\scanpage{52}

\deleted{
% 1–2
foncteurs admissibles étant maintenant les foncteurs $\Sigma$-exacts…)

% 3
\uline{Indications sur dém}

% 4–7
Dans $\mathbb P_{\bullet}$, il est évident que \uline{$\mathbb E_{\bullet}^{\bullet}$ est une $\Sigma_{\bullet}$-spéciale résolution}. De façon précise, si $\deleted{\unsure{-|\to}}\Psi_{\bullet}^{i}=\deleted{\unsure{Z^{i}(\mathbb E_{\bullet})|H^i(C^\bullet)}} k[\deleted{\unsure{\,i\,}}i]$ placé en degré $i$ on voit que l'on a des suites $\in\Sigma_{\bullet}$

% 8–10, staggered rows
\[
 \begin{aligned}
 &0\longrightarrow\Psi_{\bullet}^{0}\longrightarrow C_{\bullet}^{0}\longrightarrow\Psi_{\bullet}^{1}\longrightarrow0\\
 &\hspace{4em}0\longrightarrow\Psi_{\bullet}^{1}\longrightarrow C_{\bullet}^{1}\longrightarrow\Psi_{\bullet}^{2}\longrightarrow0\\
 &\hspace{8em}0\longrightarrow\Psi_{\bullet}^{2}\longrightarrow C_{\bullet}^{2}\longrightarrow\Psi_{\bullet}^{3}\longrightarrow\unsure{0|\cdots}
 \end{aligned}
\]

% 11–14
De plus, on voit que pour $L_{\bullet}\in\operatorname{Ob}\unsure{\overline{\mathbb P}_{\bullet}|\mathbb P_{\bullet}|\mathcal P_{\bullet}}$, $L_{\bullet}$ peut s'écrire comme extension successive d'objets de la forme $(\Psi_{\bullet}^{i})^{\unsure{\rho_i|f_i}}$ pour la filtration croissante

% 15–16
\[
 \mathrm{Fil}_i(L_{\bullet})_j=\begin{cases}L_j&\text{si }j\leq i\\0&\text{si }j>i\end{cases}
\]

% 17
on a un isom can

% 18–19
\[
 \begin{aligned}
 \deleted{\mathrm{Fil}}\ \mathrm{Gr}_i(L_{\bullet})&\simeq L_i\ \text{placé en degré }i\\
 &\underset{\mathrm{can}}{\simeq}\Psi_{\bullet}^{i}\otimes L_i\underset{\text{non can}}{\simeq}(\Psi_{\bullet}^{i})^{\rho_i}
 \end{aligned}
\]

% 20
\uline{Enfin}, les $\deleted{\unsure{\mathbb E}}C_{\bullet}^{i}$ sont $\Sigma_{\bullet}$-projectifs

% 21–22
Soit $B^{\bullet}_{\deleted{\unsure{\Sigma}}\Sigma\text{-}\mathrm{sex}}$ (resp. $B^{\bullet}_{\Sigma\text{-}\mathrm{ex}}$) la sous-catégorie pleine de $B^{\bullet}$ décrite dans a) (resp. b)).

% 23
On va définir des foncteurs

% 24–25, square
\diagram{\begin{tikzcd}[ampersand replacement=\&]
 B^{\bullet}_{\Sigma\text{-}\mathrm{sex}}\arrow[r]\arrow[d]\&\underline{\operatorname{Hom}}_{\Sigma\text{-}\mathrm{sex}}(\mathbb P_{\bullet},B)\arrow[d]\\
 B^{\bullet}_{\Sigma\text{-}\mathrm{ex}}\arrow[r]\&\underline{\operatorname{Hom}}_{\Sigma\text{-}\mathrm{ex}}(\mathbb P_{\bullet},B)
\end{tikzcd}}

% 26
(le premier induisant le deuxième,
}

% archive, bottom left
% archive sheet number: 51
